% Косой мост (полумост).
% Делитель C1-C2 создаёт среднюю точку; два ключа VT1/VT2 (MOSFET с
% индуцированным каналом) поочерёдно подключают первичную обмотку к
% половинам входного напряжения. Вторичная цепь — к ВЧ-выпрямителю.
\documentclass[border=4pt]{standalone}
\usepackage{amsmath}
\usepackage{fontspec}
\setmainfont{Calibri}
\usepackage[european, RPvoltages]{circuitikz}
\usetikzlibrary{arrows.meta, calc}
\definecolor{ink}{RGB}{31,42,51}
\definecolor{grp}{RGB}{46,110,140}
\ctikzset{bipoles/length=1.0cm}
\ctikzset{bipoles/thickness=0.9}
\ctikzset{diodes/fill=ink}
\ctikzset{inductor=american}

\begin{document}
\begin{tikzpicture}[line width=1.0pt, color=ink, font=\small]
  % --- входные зажимы (открытые)
  \draw (0.6,4.0) node[ocirc]{} node[above left]{$+$} -- (4.3,4.0);
  \node[anchor=east] at (0.4,2.0) {$U_{\text{вх}}$};
  % --- делитель C1-C2, средняя точка M
  \draw (1.5,4.0) node[circ]{} to[C, l=$C_1$] (1.5,2.2) node[circ]{}
        to[C, l=$C_2$] (1.5,0) node[circ]{};
  % --- первичная обмотка: от средней точки делителя к узлу A
  \draw (1.5,2.2) -- (2.3,2.2) to[L, l_=$w_1$] (3.7,2.2) -- (4.3,2.2)
        node[circ]{};
  % --- ключи VT1/VT2 (пути снизу вверх, затворы влево)
  \draw (4.3,2.2) node[circ]{} to[Tnigfete, name=VT1] (4.3,3.2)
        -- (4.3,4.0);
  \draw (4.3,0) node[circ]{} to[Tnigfete, name=VT2] (4.3,1.4)
        -- (4.3,2.2);
  \node[anchor=west] at (4.75,2.8) {VT1};
  \node[anchor=west] at (4.75,0.9) {VT2};
  % --- затворы: ШИМ (влево)
  \draw[grp] (VT1.G) -- ++(-0.5,0);
  \node[anchor=east, text=grp] at ($(VT1.G)+(-0.55,0)$) {ШИМ};
  \draw[grp] (VT2.G) -- ++(-0.5,0);
  \node[anchor=east, text=grp] at ($(VT2.G)+(-0.55,0)$) {ШИМ};
  % --- нижняя шина
  \draw (0.6,0) node[ocirc]{} node[below left]{$-$} -- (4.3,0);
  % --- магнитопровод (две горизонтальные линии между обмотками)
  \draw[line width=0.8pt] (2.45,1.72) -- (3.45,1.72);
  \draw[line width=0.8pt] (2.45,1.60) -- (3.45,1.60);
  \draw (2.3,1.2) to[L, l_=$w_2$] (3.7,1.2);
  % --- выводы вторички — к ВЧ-выпрямителю (развязка)
  \draw[-{Stealth}, thin, grp] (3.9,1.2) -- (4.9,1.2);
  \node[anchor=west, text=grp] at (5.0,1.2) {к ВЧ-выпрямителю};
\end{tikzpicture}
\end{document}
